\documentclass[11pt]{article}

\usepackage{nmrdoc}
\usepackage{siunitx}
\usepackage{bm}

\sisetup{per-mode=symbol}
\DeclareSIUnit{\angstrom}{\text{\AA}}
\DeclareSIUnit{\kcal}{kcal}
\emergencystretch=12em

\begin{document}
\NmrTitle{The HBond Calculator}

\section{What the calculator computes}

The HBond calculator describes a through-space tensor associated with
backbone hydrogen bonding. It starts from hydrogen-bond partners reported by
DSSP, resolves each partner to a donor backbone \(N\) atom and an acceptor
backbone \(O\) atom, and treats the \(N\cdots O\) line as an oriented source.
For every atom in the conformation, each accepted source contributes a
\(3\times 3\) geometric tensor with units \(\si{\angstrom^{-3}}\). The
observable in the background is the nuclear shielding tensor, and after
referencing it is the chemical shift, but this calculator does not output ppm
shieldings. It outputs the geometric part that can correlate with a DFT
shielding contribution after calibration.

The physical idea is that a hydrogen bond can reorganize electron density
around the donor and acceptor groups. In an applied magnetic field, that
electronic response can change the secondary magnetic field felt at a nearby
nucleus. The code does not compute the electronic response. It keeps geometry
used as a proxy for part of the angular signal: the donor-to-acceptor
direction, the source midpoint, and the vector from that midpoint to the atom
being evaluated.

\section{Where shielding enters}

For a nucleus in an applied magnetic field \(\bm B_0\), the local field is
written
\begin{equation}
  \bm B_{\mathrm{nuc}} = (\bm I-\bm\sigma)\bm B_0 .
  \label{eq:shielding}
\end{equation}
Here \(\bm I\) is the \(3\times 3\) identity matrix and \(\bm\sigma\) is the
shielding tensor. The tensor is dimensionless in physical NMR: it is the
linear coefficient connecting a small applied field to the electronic field
that opposes or augments it at the nucleus. A matrix is needed because a
protein conformation is not spherical. A field applied along \(z\), for
example, can induce a secondary field whose \(x\), \(y\), and \(z\) components
depend on how the local groups are oriented.

Equation~\eqref{eq:shielding} is a quantum response statement. The applied
field perturbs the electronic state, the perturbed electrons produce their own
magnetic field, and \(\bm\sigma\) records the linear part of that response. The
HBond calculator replaces that electronic calculation by an oriented geometric
source. It says: if DSSP reports that one peptide \(N-H\) donates to another
peptide \(C=O\), then the line from the donor \(N\) to the acceptor \(O\) is a
local axis that may organize the shielding response. The tensor form keeps
that axis as an angular object rather than reducing it to a distance count.

The solution NMR scalar most often compared across atoms is the isotropic
chemical shift, related to \(\operatorname{tr}\bm\sigma/3\) after reference
subtraction and sign conventions. This calculator keeps the full rank-2
descriptor because the anisotropic part records how the signal changes with
molecular orientation. That angular part is the content used when the
descriptor is compared with tensor-valued DFT shielding data.

\section{The tensor split used by the code}

The code stores every \(3\times 3\) tensor \(\bm X\) through the
\texttt{SphericalTensor} decomposition. This projects the matrix into three
pieces that rotate independently: a scalar, an axial vector, and a symmetric
traceless tensor. It is not an eigendecomposition. The scalar part is
\begin{equation}
  T_0 = \frac{1}{3}\operatorname{tr}\bm X ,
  \label{eq:t0}
\end{equation}
the average of the three diagonal entries. It is the isotropic component.

The antisymmetric part is stored as the three-component vector
\begin{equation}
  \bm T_1 =
  \frac{1}{2}
  \begin{pmatrix}
    X_{yz}-X_{zy}\\
    X_{zx}-X_{xz}\\
    X_{xy}-X_{yx}
  \end{pmatrix}.
  \label{eq:t1}
\end{equation}
This vector is the compact dual of the matrix part that changes sign under
transpose. It is usually not observed in standard isotropic solution NMR, but
the HBond tensor can contain it because its main term is generally asymmetric.

The remaining part is the symmetric traceless matrix
\begin{equation}
  \bm S =
  \frac{\bm X+\bm X^{\mathsf T}}{2} - T_0\bm I,
  \qquad \operatorname{tr}\bm S=0 .
  \label{eq:s}
\end{equation}
It has five independent entries. The code writes them in the real spherical
\(T_2\) basis
\begin{equation}
  \bm T_2 =
  \begin{pmatrix}
    \sqrt{2}\,S_{xy}\\
    \sqrt{2}\,S_{yz}\\
    \sqrt{3/2}\,S_{zz}\\
    \sqrt{2}\,S_{xz}\\
    (S_{xx}-S_{yy})/\sqrt{2}
  \end{pmatrix}.
  \label{eq:t2}
\end{equation}
The factors in Eq.~\eqref{eq:t2} are part of the file format. They preserve
the Frobenius norm: the squared length of this five-vector equals
\(\sum_{ab}S_{ab}^2\). HBond populates all three channels, \(T_0\), \(T_1\),
and \(T_2\), by decomposing the accumulated full tensor for each atom.

\section{The method implemented}

Hydrogen-bond sources come from \texttt{DsspResult}. Here,
\texttt{DsspResult} calls \texttt{libdssp}, which implements the DSSP
Kabsch--Sander hydrogen-bond assignment. \texttt{HBondResult} does not
recompute that energy, inspect the stored energy values, or apply a DSSP
energy threshold. It uses every partner residue index that
\texttt{DsspResult} reports, then applies its own geometric filters.

For each residue \(i\), the code reads up to two acceptor partners and two
donor partners. On the donor side, residue \(i\)'s backbone \(N\) donates to a
partner residue's backbone \(O\). On the acceptor side, a partner residue's
backbone \(N\) donates to residue \(i\)'s backbone \(O\). A candidate source is
discarded if either atom is missing, if the donor and acceptor residues are
less than \(\num{2}\) residues apart in sequence, if the pair has already been
seen, if the \(N\cdots O\) distance is below \(\SI{0.1}{\angstrom}\), or if
that distance is above \(\SI{50.0}{\angstrom}\). The last two values are
\texttt{singularity\_guard\_distance} and
\texttt{hbond\_dipolar\_max\_distance} in
\texttt{data/calculator\_params.toml}. No local donor--hydrogen--acceptor
angle test appears in this calculator.

An exact way to hold the model is to replace each accepted hydrogen bond by a
short oriented segment from donor \(N\) to acceptor \(O\), then keep its
midpoint and direction. Those are the objects the code stores. If the donor and
acceptor positions are \(\bm n_k\) and \(\bm o_k\), then
\begin{equation}
  \bm m_k = \frac{\bm n_k+\bm o_k}{2},\qquad
  L_k = |\bm o_k-\bm n_k|,\qquad
  \hat{\bm h}_k = \frac{\bm o_k-\bm n_k}{L_k}.
  \label{eq:source}
\end{equation}
\(\bm m_k\) is the source point in \(\si{\angstrom}\), \(L_k\) is the
\(N\cdots O\) length in \(\si{\angstrom}\), and \(\hat{\bm h}_k\) is a
unitless direction from donor to acceptor.

For a target atom at \(\bm r_i\), the code forms
\begin{equation}
  \bm d_{ik} = \bm r_i-\bm m_k,\qquad
  r_{ik}=|\bm d_{ik}|,\qquad
  \hat{\bm d}_{ik}=\frac{\bm d_{ik}}{r_{ik}},\qquad
  c_{ik}=\hat{\bm d}_{ik}\cdot\hat{\bm h}_k .
  \label{eq:obs}
\end{equation}
Here \(r_{ik}\) is the atom-to-midpoint distance in \(\si{\angstrom}\), and
\(c_{ik}\) is the cosine of the angle between the source direction and the
midpoint-to-atom direction. The dyadic product \(u_a v_b\) is a \(3\times 3\)
table: it takes a component along \(\bm v\) and returns a vector along
\(\bm u\). The HBond contribution is
\begin{equation}
  M^{(ik)}_{ab} =
  \frac{
    9c_{ik}\,\hat d_{ik,a}\hat h_{k,b}
    -3\hat h_{k,a}\hat h_{k,b}
    -\left(3\hat d_{ik,a}\hat d_{ik,b}-\delta_{ab}\right)}
  {r_{ik}^{3}} .
  \label{eq:hbond}
\end{equation}
The Cartesian indices \(a,b\) run over \(x,y,z\), and \(\delta_{ab}\) is one
when \(a=b\) and zero otherwise. The numerator is dimensionless, so
\(M^{(ik)}\) has units \(\si{\angstrom^{-3}}\). The first term uses both the
viewing direction \(\hat{\bm d}\) and the hydrogen-bond direction
\(\hat{\bm h}\), which is why the tensor can be asymmetric. The last term is a
traceless line-of-sight tensor. Taking the trace gives
\begin{equation}
  \frac{1}{3}\operatorname{tr}\bm M^{(ik)}
  =
  f_{ik}
  =
  \frac{3c_{ik}^{2}-1}{r_{ik}^{3}} .
  \label{eq:f}
\end{equation}
The scalar \(f_{ik}\) is accumulated as a separate feature over the same
accepted atom--source pairs.

Before adding Eq.~\eqref{eq:hbond}, the code applies three atom--source
filters. First, \(r_{ik}\) is accepted at or above \(\SI{0.1}{\angstrom}\),
preventing the \(r^{-3}\) expression from being evaluated at a singular point.
Second, the target atom may not be the donor \(N\) or acceptor \(O\) of that
source. Third, the target is accepted outside the near field of the finite
\(N\cdots O\) segment:
\begin{equation}
  r_{ik} > \alpha L_k,\qquad \alpha=0.5 .
  \label{eq:nearfield}
\end{equation}
The parameter \(\alpha\) is \texttt{near\_field\_exclusion\_ratio}. The point
source description is not used inside half the segment length. The calculator
does not use a spatial search cutoff for atom--source evaluation; after DSSP
has supplied the source list, all sources are considered and the \(r^{-3}\)
decay sets their size.

For each atom, the stored tensor is
\begin{equation}
  \bm M_i = \sum_{k\in \mathcal A_i}\bm M^{(ik)},
  \label{eq:sum}
\end{equation}
where \(\mathcal A_i\) is the set of resolved hydrogen bonds that pass the
filters for atom \(i\). The nearest accepted source is tracked by \(r_{ik}\).
The code also counts accepted sources with \(r_{ik}<\SI{3.5}{\angstrom}\), the
configured \texttt{hbond\_counting\_radius}. The nearest-distance and count
features are geometric diagnostics; they do not change the tensor sum.
The \texttt{SampleKernelAt} path evaluates the same sum at an arbitrary point
using the stored source list. Because that point has no atom index, it applies
the singularity and near-field checks but not the donor/acceptor endpoint
exclusion.

The main approximations are visible in these equations. Each hydrogen bond is
reduced to one midpoint and one \(N\to O\) axis. All accepted sources have unit
weight; the code does not multiply by DSSP energy, hydrogen-bond length
dependence, partial charge, susceptibility, or an empirical ppm coefficient. It
uses backbone DSSP partners, so side-chain and water hydrogen bonds are outside
this calculator. It uses \(N\to O\), not \(H\to O\), as the tensor axis after
DSSP has identified the donor and acceptor.

\section{External inputs and output}

HBondResult consumes atom positions and residue topology from the current
protein conformation, plus DSSP partner data from \texttt{DsspResult}. DSSP is
the external library boundary: \texttt{DsspResult} calls \texttt{libdssp}, then
maps its secondary-structure and backbone hydrogen-bond partner slots back onto
the conformation. \texttt{SpatialIndexResult} is declared as a dependency, but
this calculator iterates over the resolved hydrogen bonds directly rather than
querying a neighbor list.

MOPAC charges and bond orders, AIMNet2 charges or embeddings, ff14SB partial
charges, and APBS electric fields are not inputs to this calculator. They feed
other environment and electrostatic descriptors in the pipeline. Here the
written files are \texttt{hbond\_shielding.npy}, an \(N\times 9\) array in the
\([T_0,T_1,T_2]\) layout of Eqs.~\eqref{eq:t0}--\eqref{eq:t2}, and
\texttt{hbond\_scalars.npy}, an \(N\times 4\) array containing nearest
atom-to-midpoint distance, its \(1/r^3\), the count inside
\(\SI{3.5}{\angstrom}\), and \(\sum f_{ik}\). If no source passes the filters
for an atom, these scalar columns and the tensor remain at their zero initial
values. Despite the file name, the tensor array is a geometric quantity in
\(\si{\angstrom^{-3}}\). Calibration against DFT reference tensors is what maps
it to a shielding contribution.

\end{document}
